\par\Needspace{6\baselineskip}\medskip
\hypertarget{native-218-article159}{}\pdfbookmark[2]{159.}{article-159}
\begin{center}159.\end{center}
Für das \emph{dritte} Beispiel wählen wir die folgenden Beobachtungen der Ceres. Die
erste stellte der namhafte \textsc{Olbers} in Bremen an, die zweite der namhafte \textsc{Harding} in Göttingen und die dritte der namhafte \textsc{Bessel} in Lilienthal:
\begin{center}\begingroup\footnotesize\setlength{\tabcolsep}{3pt}\renewcommand{\arraystretch}{1.3}
\begin{tabular}{|c|c|c|}
\hline
\shortstack{Mittlere Ortszeit\\der Beobachtung} & Rektasz. & Nördl. Dekl. \\
\hline
1805 Sept. 5. $13^{\mathrm h}$ $8^{\mathrm m}$ $54^{\mathrm s}$ & $95^0\,59'\,25''$ & $22^0\,21'\,25''$ \\
1806 Jan. 17. $10$ $58$ $51$ & $101\quad18\quad40{,}6$ & $30\quad21\quad22{,}3$ \\
1806 Mai 23. $10$ $23$ $53$ & $121\quad56\quad7$ & $28\quad2\quad45$ \\
\hline
\end{tabular}
\endgroup\end{center}

Die beiden vorangegangenen Beispiele haben bereits ausführlich gezeigt, wie sich
Parallaxe und Aberration berücksichtigen lassen, wenn die Abstände von der Erde
als völlig unbekannt gelten. In diesem dritten Beispiel ersparen wir uns daher den
Mehraufwand und entnehmen der gelehrten Korrespondenz des namhaften
de \textsc{Zach} (Bd. XI S. 284) Näherungswerte für die Abstände, um die Beobachtungen von den Auswirkungen der Parallaxe und Aberration zu befreien. Diese Abstände
und die daraus abgeleiteten Reduktionen gibt die folgende Tabelle an:
\begin{center}\begingroup\footnotesize\setlength{\tabcolsep}{3pt}\renewcommand{\arraystretch}{1.3}
\begin{tabular}{l|c|c|c}
Abstand der Ceres von der Erde & $2{,}899$ & $1{,}638$ & $2{,}964$ \\
\shortstack[l]{Laufzeit des Lichts\\bis zur Erde} & $23^{\mathrm m}49^{\mathrm s}$ & $13^{\mathrm m}28^{\mathrm s}$ & $24^{\mathrm m}21^{\mathrm s}$ \\
reduzierter Beobachtungszeitpunkt & $12^{\mathrm h}45^{\mathrm m}5^{\mathrm s}$ & $10^{\mathrm h}45^{\mathrm m}23^{\mathrm s}$ & $9^{\mathrm h}59^{\mathrm m}32^{\mathrm s}$ \\
Sternzeit in Grad & $355^0\,55'$ & $97^0\,59'$ & $210^0\,41'$ \\
Parallaxe in Rektaszension & $+1''{,}90$ & $+0''{,}22$ & $-1''{,}97$ \\
Parallaxe in Deklination & $-2{,}08$ & $-1{,}90$ & $-2{,}04$ \\
\end{tabular}
\endgroup\end{center}

\SourceLeaf{219}Nach der Korrektur für Parallaxe und Aberration und der Reduktion der Zeiten
auf den Pariser Meridian lauten die Daten der Aufgabe somit:
\begin{center}\begingroup\footnotesize\setlength{\tabcolsep}{3pt}\renewcommand{\arraystretch}{1.3}
\begin{tabular}{|c|c|c|}
\hline
 & Rektasz. & Deklination \\
\hline
1805 Sept. 5. $12^{\mathrm h}$ $19^{\mathrm m}$ $14^{\mathrm s}$ & $95^0\,59'\,23''{,}10$ & $22^0\,21'\,27''{,}08$ \\
1806 Jan. 17. $10$ $15$ $2$ & $101\quad18\quad40{,}38$ & $30\quad21\quad24{,}20$ \\
1806 Mai 23. $9$ $33$ $18$ & $121\quad56\quad8{,}97$ & $28\quad2\quad47{,}04$ \\
\hline
\end{tabular}
\endgroup\end{center}

Aus diesen Rektaszensionen und Deklinationen wurden die Längen und Breiten
mit den Werten $23^0\,27'\,55''{,}90$, $23^0\,27'\,54''{,}59$,
$23^0\,27'\,53''{,}27$ für die Schiefe der Ekliptik abgeleitet. Anschließend wurden die Längen um die Nutation korrigiert, die jeweils
$+17''{,}31$, $+17''{,}88$, $+18''{,}00$ betrug, und auf den Jahresanfang 1806 reduziert, wobei die Präzession $+15''{,}98$, $-2''{,}39$, $-19''{,}68$ berücksichtigt wurde. Schließlich entnahmen wir
den Tafeln die Orte der Sonne für die reduzierten Zeiten. Bei den Längen blieb
die Nutation unberücksichtigt, während die Präzession ebenso wie bei den Längen der Ceres hinzugerechnet wurde. Die Breite der Sonne
wurde ganz vernachlässigt. So ergaben sich die folgenden Zahlen für die Rechnung:
\begin{center}\begingroup\footnotesize\setlength{\tabcolsep}{3pt}\renewcommand{\arraystretch}{1.3}
\begin{tabular}{l|c|c|c}
Zeit 1805 Sept. & $5{,}51336$ & $139{,}42711$ & $265{,}39813$ \\
$a$, $a'$, $a''$ & $95^0\,32'\,18''{,}56$ & $99^0\,49'\,5''{,}87$ & $118^0\,5'\,28''{,}85$ \\
$\beta$, $\beta'$, $\beta''$ & $-0\quad59\quad34{,}06$ & $+7\quad16\quad36{,}80$ & $+7\quad38\quad49{,}39$ \\
$l$, $l'$, $l''$ & $342\quad54\quad56{,}00$ & $117\quad12\quad43{,}25$ & $241\quad58\quad50{,}71$ \\
$\log R$, $\log R'$, $\log R''$ & $0{,}003\,1514$ & $9{,}992\,9861$ & $0{,}005\,6974$ \\
\end{tabular}
\endgroup\end{center}

Die in Art. 136---143 erläuterten vorbereitenden Rechnungen ergeben nun Folgendes:
\begin{center}\begingroup\footnotesize\setlength{\tabcolsep}{3pt}\renewcommand{\arraystretch}{1.3}
\begin{tabular}{l|c|c|c}
$\gamma$, $\gamma'$, $\gamma''$ & $358^0\,55'\,28''{,}09$ & $156^0\,52'\,11''{,}49$ & $170^0\,48'\,44''{,}79$ \\
$\delta$, $\delta'$, $\delta''$ & $112\quad37\quad9{,}66$ & $18\quad48\quad39{,}81$ & $123\quad32\quad52{,}13$ \\
$A'D$, $AD'$, $AD''$ & $15\quad32\quad41{,}40$ & $252\quad42\quad19{,}14$ & $136\quad2\quad22{,}38$ \\
$A''D$, $A''D'$, $A'D''$ & $138\quad45\quad4{,}60$ & $6\quad26\quad41{,}10$ & $358\quad5\quad57{,}00$ \\
$\epsilon$, $\epsilon'$, $\epsilon''$ & $29\quad18\quad8{,}21$ & $170\quad32\quad59{,}08$ & $156\quad6\quad25{,}25$ \\
\end{tabular}
\endgroup\end{center}
\SourceLeaf{220}
\begin{center}\begingroup\footnotesize\setlength{\tabcolsep}{3pt}\renewcommand{\arraystretch}{1.3}
\begin{tabular}{ll}
$\sigma=8^0\,52'\,4''{,}05$ & \\
$\log a=0{,}184\,0193_n$, & $a=-1{,}527\,6340$\\
$\log b=0{,}004\,0987$ & \\
$\log\varkappa=0{,}161\,1012$ & \\
$\log\lambda=9{,}916\,4090_n$ &
\end{tabular}
\par\smallskip
\begin{tabular}{r@{\;=\;}l}
$\log c$ & $2{,}006\,6735$\\
$d$ & $117{,}50873$\\
$\log e$ & $0{,}856\,8244$\\
$\log\varkappa''$ & $9{,}977\,0819_n$\\
$\log\lambda''$ & $9{,}732\,0127_n$.
\end{tabular}
\endgroup\end{center}

Der zeitliche Abstand zwischen der ersten und der zweiten Beobachtung beträgt
$133{,}91375$ Tage, zwischen der zweiten und der dritten $125{,}97102$ Tage. Daraus ergibt sich $\log\theta = 0{,}335\,8520$,
$\log\theta'' = 0{,}362\,4066$, $\log\dfrac{\theta''}{\theta} = 0{,}026\,5546$, $\log\theta\theta'' = 0{,}698\,2586$. Die wichtigsten Größen der ersten drei nacheinander
aufgestellten Hypothesen geben wir in der folgenden Übersicht an:
\begin{center}\begingroup\footnotesize\setlength{\tabcolsep}{3pt}\renewcommand{\arraystretch}{1.3}
\begin{tabular}{l|c|c|c}
 & I & II & III \\
\hline
$\log P = x$ & $0{,}026\,5546$ & $0{,}025\,6968$ & $0{,}025\,6275$ \\
$\log Q = y$ & $0{,}698\,2586$ & $0{,}739\,0190$ & $0{,}748\,1055$ \\
$\omega + \sigma$ & $7^0\,15'\,13''{,}523$ & $7^0\,14'\,47''{,}199$ & $7^0\,14'\,45''{,}071$ \\
$\log Qc\sin\omega$ & $1{,}154\,6650_n$ & $1{,}197\,3880_n$ & $1{,}206\,6327_n$ \\
$z$ & $7^0\,3'\,59''{,}018$ & $7^0\,2'\,32''{,}938$ & $7^0\,2'\,16''{,}900$ \\
$\log r'$ & $0{,}411\,4726$ & $0{,}412\,9371$ & $0{,}413\,2107$ \\
$\zeta$ & $160^0\,10'\,46''{,}74$ & $160^0\,20'\,7''{,}82$ & $160^0\,22'\,9''{,}42$ \\
$\zeta''$ & $262\quad6\quad1{,}03$ & $262\quad12\quad18{,}26$ & $262\quad14\quad19{,}49$ \\
$\log r$ & $0{,}432\,3934$ & $0{,}429\,1773$ & $0{,}428\,4841$ \\
$\log r''$ & $0{,}409\,4712$ & $0{,}407\,1975$ & $0{,}406\,4697$ \\
$\tfrac12(u''+u)$ & $262^0\,55'\,23''{,}22$ & $262^0\,57'\,6''{,}83$ & $262^0\,57'\,31''{,}17$ \\
$\tfrac12(u''-u)$ & $273\quad28\quad50{,}95$ & $273\quad29\quad15{,}06$ & $273\quad29\quad19{,}56$ \\
$2f'$ & $62\quad34\quad28{,}40$ & $62\quad49\quad56{,}50$ & $62\quad53\quad57{,}06$ \\
$2f$ & $31\quad8\quad30{,}03$ & $31\quad15\quad59{,}09$ & $31\quad18\quad13{,}83$ \\
$2f''$ & $31\quad25\quad58{,}45$ & $31\quad33\quad57{,}32$ & $31\quad35\quad43{,}32$ \\
$\log \eta$ & $0{,}020\,2496$ & $0{,}020\,3158$ & $0{,}020\,3494$ \\
$\log \eta''$ & $0{,}021\,1074$ & $0{,}021\,2429$ & $0{,}021\,2751$ \\
$\log P'$ & $0{,}025\,6968$ & $0{,}025\,6275$ & $0{,}025\,6289$ \\
$\log Q'$ & $0{,}739\,0190$ & $0{,}748\,1055$ & $0{,}750\,2337$ \\
$X$ & $-0{,}000\,8578$ & $-0{,}000\,0693$ & $+0{,}000\,0014$ \\
$Y$ & $+0{,}040\,7604$ & $+0{,}009\,0865$ & $+0{,}002\,1282$ \\
\end{tabular}
\endgroup\end{center}
\SourceLeaf{221}

Wir bezeichnen nun die drei Werte von $X$ mit $A$, $A'$, $A''$ und die drei Werte von $Y$
mit $B$, $B'$, $B''$. Die Quotienten, die sich bei der Division von $A'B''-A''B'$, $A''B-AB''$,
$AB'-A'B$ durch ihre Summe ergeben, bezeichnen wir mit $k$, $k'$, $k''$ in dieser Reihenfolge, sodass $k+k'+k''=1$ gilt. Schließlich bezeichnen wir die Werte von $\log P'$ und $\log Q'$ in der dritten
Hypothese mit $M$ und $N$ (dies wären die neuen Werte von $x$, $y$, wenn man die vierte
Hypothese aus der dritten ebenso ableiten wollte wie die dritte aus
der zweiten). Aus den Formeln in Art. 120 ergibt sich leicht, dass der korrigierte Wert von $x$
$=M-k(A'+A'')-k'A''$ beträgt und dass für den anderen korrigierten Wert $y = N-k(B'+B'')-k'B''$ gilt.
Die Rechnung ergibt für den ersten Wert $=0{,}025\,6331$, für den zweiten $=0{,}750\,9143$. Auf diese korrigierten Werte gründen wir nun die \emph{vierte Hypothese}, deren wichtigste Größen die folgenden sind:
\begin{center}\begingroup\footnotesize\setlength{\tabcolsep}{3pt}\renewcommand{\arraystretch}{1.3}
\begin{tabular}{r|r}
$\omega + \sigma=7^0\,14'\,45''{,}247$ & $\log r''=0{,}406\,2033$\\
$\log Qc\sin\omega=1{,}209\,4284_n$ & $\tfrac12(u''+u)=262^0\,57'\,38''{,}78$\\
$z=7^0\,2'\,12''{,}736$ & $\tfrac12(u''-u)=273\quad29\quad20{,}73$\\
$\log r'=0{,}413\,2817$ & $2f'=62\quad55\quad16{,}64$\\
$\zeta=160^0\,22'\,45''{,}38$ & $2f=31\quad19\quad1{,}49$\\
$\zeta''=262\quad15\quad3{,}90$ & $2f''=31\quad36\quad15{,}20$\\
$\log r=0{,}428\,2792$ & \\
\end{tabular}
\endgroup\end{center}

Zwischen $2f'$ und $2f+2f''$ ergibt sich ein Unterschied von $0''{,}05$, den wir so verteilen, dass
wir $2f=31^0\,19'\,1''{,}47$, $2f''=31^0\,36'\,15''{,}17$ setzen. Bestimmt man nun die Elemente selbst aus den beiden
äußersten Orten, so ergeben sich die folgenden Zahlen:
\begin{center}\begingroup\footnotesize\setlength{\tabcolsep}{3pt}\renewcommand{\arraystretch}{1.3}
\begin{tabular}{ll}
Wahre Anomalie am ersten Ort\quad .\,.\,.\,.\,.\,.\,.\,.\,.\,.\,.\,.\,.\,.\,.\,. & $289^0\,7'\,39''{,}75$\\
Wahre Anomalie am dritten Ort\quad .\,.\,.\,.\,.\,.\,.\,.\,.\,.\,.\,.\,.\,.\,.\,. & $352\quad2\quad56{,}39$\\
Mittlere Anomalie am ersten Ort\quad .\,.\,.\,.\,.\,.\,.\,.\,.\,.\,.\,.\,.\,.\,. & $297\quad41\quad35{,}65$\\
Mittlere Anomalie am dritten Ort\quad .\,.\,.\,.\,.\,.\,.\,.\,.\,.\,.\,.\,.\,.\,. & $353\quad15\quad22{,}49$\\
Mittlere tägliche siderische Bewegung\quad .\,.\,.\,.\,.\,.\,.\,.\,.\,.\,.\,.\,.\,.\,.\,.\,. & $769''{,}6750$\\
Mittlere Anomalie zum Jahresanfang 1806\quad .\,.\,.\,.\,.\,.\,.\,.\,.\,. & $322^0\,35'\,52''{,}51$\\
Winkel $\varphi$\quad .\,.\,.\,.\,.\,.\,.\,.\,.\,.\,.\,.\,.\,.\,.\,.\,.\,.\,.\,.\,.\,.\,.\,.\,.\,.\,.\,.\,.\,.\,.\,.\,. & $4\quad37\quad57{,}78$\\
Logarithmus der großen Halbachse\quad .\,.\,.\,.\,.\,.\,.\,.\,.\,.\,.\,.\,.\,.\,.\,.\,. & $0{,}442\,4661$.
\end{tabular}
\endgroup\end{center}

Berechnet man mit diesen Elementen den heliozentrischen Ort zum Zeitpunkt der mittleren Beobachtung, so erhält man die mittlere Anomalie $326^0\,19'\,25''{,}72$, den Logarithmus des
Radiusvektors $0{,}413\,2825$ und die wahre Anomalie $320^0\,43'\,54''{,}87$. Letztere müsste sich von der \SourceLeaf{222} wahren Anomalie am ersten Ort um $2f''$ oder von der wahren Anomalie
am dritten Ort um $2f$ unterscheiden und somit $=320^{0}\,43'\,54''{,}92$ betragen. Ebenso müsste
der Logarithmus des Radiusvektors $=0{,}413\,2817$ betragen. Der Unterschied von $0''{,}05$ in der wahren Anomalie
und der von acht Einheiten in diesem Logarithmus sind als unerheblich anzusehen.

Führte man die vierte Hypothese ebenso zu Ende wie die drei vorangegangenen,
so ergäbe sich $X=0$, $Y=-0{,}000\,0168$. Daraus ließen sich die folgenden korrigierten Werte für $x$, $y$
ableiten
\[
\begin{aligned}x &= \log P = 0{,}025\,6331 \\ &\quad\text{(wie in der vierten Hypothese)}
\\
y &= \log Q = 0{,}750\,8973.
\end{aligned}
\]
\noindent Gründete man auf diese Werte eine fünfte Hypothese, so erreichte die Lösung die höchste
Genauigkeit, die die Tafeln zulassen. Die daraus gewonnenen Elemente würden sich aber
kaum merklich von denen der vierten Hypothese unterscheiden.

\Needspace{19\baselineskip}
Um die vollständigen Elemente zu erhalten, braucht man nur noch die Lage der Bahnebene zu
berechnen. Nach den Regeln in Art. 149 ergibt sich
\begin{center}\begingroup\footnotesize\setlength{\tabcolsep}{3pt}\renewcommand{\arraystretch}{1.3}
\begin{tabular}{ll}\multicolumn{2}{c}{aus dem ersten Ort}\\
$g$\quad .\,.\,.\,.\,.\,.\,.\,.\,.\,.\,.\,.\,.\,.\,.\,.\,.\,.\,.\,.\,.\,.\,.\,.\,.\,.\,.\,. & $354^0\,9'\,44''{,}22$\\
$h$\quad .\,.\,.\,.\,.\,.\,.\,.\,.\,.\,.\,.\,.\,.\,.\,.\,.\,.\,.\,.\,.\,.\,.\,.\,.\,.\,.\,. & $261\quad56\quad6{,}94$\\
$i$\quad .\,.\,.\,.\,.\,.\,.\,.\,.\,.\,.\,.\,.\,.\,.\,.\,.\,.\,.\,.\,.\,.\,.\,.\,.\,.\,.\,. & $10\quad37\quad33{,}02$\\
$\AscendingNode$\quad .\,.\,.\,.\,.\,.\,.\,.\,.\,.\,.\,.\,.\,.\,.\,.\,.\,.\,.\,.\,.\,.\,.\,.\,.\,.\,. & $80\quad58\quad49{,}06$\\
Abstand des Perihels vom aufsteigenden Knoten & $65\quad2\quad4{,}47$\\
Länge des Perihels\quad .\,.\,.\,.\,.\,.\,.\,.\,.\,.\,.\,.\,. & $146\quad0\quad53{,}53$\\
\noalign{\smallskip}\multicolumn{2}{c}{aus dem dritten Ort}\\
$g''$\quad .\,.\,.\,.\,. & $57^0\,5'\,0''{,}91$\\
$h''$\quad .\,.\,.\,.\,. & $161\quad0\quad1{,}61$\\
.\,.\,.\,.\,.\,.\,. & $10\quad37\quad33{,}00$\\
.\,.\,.\,.\,.\,.\,. & $80\quad58\quad49{,}10$\\
.\,.\,.\,.\,.\,.\,. & $65\quad2\quad4{,}52$\\
.\,.\,.\,.\,.\,.\,. & $146\quad0\quad53{,}62$.
\end{tabular}
\endgroup\end{center}
\noindent Als Mittelwerte setzen wir daher $i=10^{0}\,37'\,33''{,}01$, $\AscendingNode=80^{0}\,58'\,49''{,}08$,
Länge des Perihels $=146^{0}\,0'\,53''{,}57$. Schließlich beträgt die mittlere Länge zum Jahresanfang
1806 $=108^{0}\,36'\,46''{,}08$.
